% SMPdesign/beyond.tex
% mainfile: ../perfbook.tex
% SPDX-License-Identifier: CC-BY-SA-3.0

\section{超越分区}
\label{sec:SMPdesign:Beyond Partitioning}
\OriginallyPublished{sec:SMPdesign:Beyond Partitioning}{Retrofitted Parallelism Considered Grossly Sub-Optimal}{4\textsuperscript{th} USENIX Workshop on Hot Topics on Parallelism}{PaulEMcKenney2012HOTPARsuboptimal}
%
\epigraph{如果你有充足的弹药，高瞄准是完全可以的。}
	 {Hawley R. Everhart}

本章讨论了如何运用数据分区来设计
既简单又能线性扩展的并行程序。
\Cref{sec:SMPdesign:Data Ownership}提到了数据复制的可能性，
这将在\cref{sec:defer:Read-Copy Update (RCU)}中发挥重大作用。

运用分区和复制的主要目标是达到线性的加速比，
换句话说，确保所需的总工作量不会随着CPU或线程数量
的增加而显著增长。
通过分区和/或复制可以解决的问题，
如果能产生线性加速比，则称为\emph{\IX{embarrassingly parallel}}（令人尴尬的并行）。
但我们还能做得更好吗？

为了回答这个问题，让我们来研究迷宫问题的求解。
千年以来，迷宫一直是个令人着迷的研究对象~\cite{WikipediaLabyrinth}，
所以用计算机来生成和求解迷宫毫不令人惊讶，
包括生物计算机~\cite{AndrewAdamatzky2011SlimeMold}、
GPGPU~\cite{ChristerEricson2008GPUMaze}，甚至
离散硬件~\cite{MIT:TRMag:MemristorMazes}。
迷宫的并行求解有时被用作大学的课堂项目~\cite{ETHZurich:FS2011maze,UMD:CMSC433maze}，
也被用作展示并行编程框架优势的工具~\cite{RonFosner2010maze}。

常见的解法是使用并行工作队列算法
（PWQ）~\cite{ETHZurich:FS2011maze,RonFosner2010maze}。
本节比较PWQ、顺序算法（SEQ）以及
另一种并行算法，
这些方法都能求解任何随机生成的方形迷宫。
\Cref{sec:SMPdesign:Work-Queue Parallel Maze Solver}讨论PWQ，
\cref{sec:SMPdesign:Alternative Parallel Maze Solver}讨论一种替代的
并行算法，
\cref{sec:SMPdesign:Performance Comparison I}分析其异常的性能，
\cref{sec:SMPdesign:Alternative Sequential Maze Solver}从替代并行算法
推导出一种改进的顺序算法，
\cref{sec:SMPdesign:Performance Comparison II}进行进一步的性能比较，
最后\cref{sec:SMPdesign:Future Directions and Conclusions}
给出未来方向和总结性评论。

\subsection{工作队列并行迷宫求解器}
\label{sec:SMPdesign:Work-Queue Parallel Maze Solver}

PWQ基于SEQ，SEQ的伪代码如
\cref{lst:SMPdesign:SEQ Pseudocode}
（\path{maze_seq.c}的伪代码）所示。
迷宫由一个二维单元格数组和
一个名为\co{->visited}的基于线性数组的工作队列表示。

\begin{listing}
\begin{fcvlabel}[ln:SMPdesign:SEQ Pseudocode]
\begin{VerbatimL}[commandchars=\\\@\$]
int maze_solve(maze *mp, cell sc, cell ec)
{
	cell c = sc;
	cell n;
	int vi = 0;

	maze_try_visit_cell(mp, c, c, &n, 1);		\lnlbl@initcell$
	for (;;) {					\lnlbl@loop:b$
		while (!maze_find_any_next_cell(mp, c, &n)) {\lnlbl@loop2:b$
			if (++vi >= mp->vi)		\lnlbl@ifge$
				return 0;
			c = mp->visited[vi].c;
		}					\lnlbl@loop2:e$
		do {					\lnlbl@loop3:b$
			if (n == ec) {
				return 1;
			}
			c = n;
		} while (maze_find_any_next_cell(mp, c, &n));\lnlbl@loop3:e$
		c = mp->visited[vi].c;			\lnlbl@finalize$
	}						\lnlbl@loop:e$
}
\end{VerbatimL}
\end{fcvlabel}
\caption{SEQ伪代码}
\label{lst:SMPdesign:SEQ Pseudocode}
\end{listing}

\begin{fcvref}[ln:SMPdesign:SEQ Pseudocode]
\Clnref{initcell}访问初始单元格，跨越
\clnrefrange{loop:b}{loop:e}的循环的每次迭代遍历以一个单元格为首的通道。
跨越\clnrefrange{loop2:b}{loop2:e}的循环扫描\co{->visited[]}数组，
寻找具有未访问邻居的已访问单元格，
而跨越\clnrefrange{loop3:b}{loop3:e}的循环遍历
以该邻居为首的子迷宫的一个分叉。
\Clnref{finalize}为外循环的下一次遍历进行初始化。
\end{fcvref}

\begin{listing}
\begin{fcvlabel}[ln:SMPdesign:SEQ Helper Pseudocode]
\begin{VerbatimL}[commandchars=\\\@\$]
int maze_try_visit_cell(struct maze *mp, cell c, cell t, \lnlbl@try:b$
                        cell *n, int d)
{
	if (!maze_cells_connected(mp, c, t) ||	\lnlbl@try:chk:adj$
	    (*celladdr(mp, t) & VISITED))	\lnlbl@try:chk:not:visited$
		return 0;			\lnlbl@try:ret:failure$
	*n = t;					\lnlbl@try:nextcell$
	mp->visited[mp->vi] = t;		\lnlbl@try:recordnext$
	mp->vi++;				\lnlbl@try:next:visited$
	*celladdr(mp, t) |= VISITED | d;	\lnlbl@try:mark:visited$
	return 1;				\lnlbl@try:ret:success$
}						\lnlbl@try:e$

int maze_find_any_next_cell(struct maze *mp, cell c, \lnlbl@find:b$
                            cell *n)
{
	int d = (*celladdr(mp, c) & DISTANCE) + 1;	\lnlbl@find:curplus1$

	if (maze_try_visit_cell(mp, c, prevcol(c), n, d))\lnlbl@find:chk:prevcol$
		return 1;				\lnlbl@find:ret:prevcol$
	if (maze_try_visit_cell(mp, c, nextcol(c), n, d))\lnlbl@find:chk:nextcol$
		return 1;				\lnlbl@find:ret:nextcol$
	if (maze_try_visit_cell(mp, c, prevrow(c), n, d))\lnlbl@find:chk:prevrow$
		return 1;				\lnlbl@find:ret:prevrow$
	if (maze_try_visit_cell(mp, c, nextrow(c), n, d))\lnlbl@find:chk:nextrow$
		return 1;				\lnlbl@find:ret:nextrow$
	return 0;					\lnlbl@find:ret:false$
}							\lnlbl@find:e$
\end{VerbatimL}
\end{fcvlabel}
\caption{SEQ辅助函数伪代码}
\label{lst:SMPdesign:SEQ Helper Pseudocode}
\end{listing}

\begin{fcvref}[ln:SMPdesign:SEQ Helper Pseudocode:try]
\co{maze_try_visit_cell()}的伪代码如
\cref{lst:SMPdesign:SEQ Helper Pseudocode}
（\path{maze.c}）的\clnrefrange{b}{e}所示。
\Clnref{chk:adj}检查单元格\co{c}和\co{t}是否
相邻且连通，
而\clnref{chk:not:visited}检查单元格\co{t}是否
尚未被访问。
\co{celladdr()}函数返回指定单元格的地址。
如果任一检查失败，\clnref{ret:failure}返回失败。
\Clnref{nextcell}指示下一个单元格，
\clnref{recordnext}将此单元格记录在
\co{->visited[]}数组的下一个槽中，
\clnref{next:visited}指示该槽现在已满，
\clnref{mark:visited}将此单元格标记为已访问并记录
距迷宫起点的距离。
\Clnref{ret:success}然后返回成功。
\end{fcvref}

\begin{fcvref}[ln:SMPdesign:SEQ Helper Pseudocode:find]
\co{maze_find_any_next_cell()}的伪代码如
\cref{lst:SMPdesign:SEQ Helper Pseudocode}
（\path{maze.c}）的\clnrefrange{b}{e}所示。
\Clnref{curplus1}获取当前单元格的距离加1，
而\clnref{chk:prevcol,chk:nextcol,chk:prevrow,chk:nextrow}
检查每个方向的单元格，
\clnref{ret:prevcol,ret:nextcol,ret:prevrow,ret:nextrow}
如果对应单元格是候选的下一个单元格则返回true。
\co{prevcol()}、\co{nextcol()}、\co{prevrow()}和\co{nextrow()}
各自执行指定的数组索引转换操作。
如果没有候选单元格，\clnref{ret:false}返回false。
\end{fcvref}

\begin{figure}
\centering
\resizebox{1.2in}{!}{\includegraphics{SMPdesign/MazeNumberPath}}
\caption{单元格编号求解跟踪}
\label{fig:SMPdesign:Cell-Number Solution Tracking}
\end{figure}

路径通过计算从起点开始的单元格数量记录在迷宫中，
如\cref{fig:SMPdesign:Cell-Number Solution Tracking}所示，
其中起始单元格在左上角，终止单元格在右下角。
从终止单元格开始，沿着连续递减的单元格编号遍历即可得到解。

\begin{figure}
\centering
\resizebox{2.2in}{!}{\includegraphics{SMPdesign/500-ms_seq_fg-cdf}}
\caption{SEQ和PWQ求解时间的CDF}
\label{fig:SMPdesign:CDF of Solution Times For SEQ and PWQ}
\end{figure}

并行工作队列求解器是对
\cref{lst:SMPdesign:SEQ Pseudocode,lst:SMPdesign:SEQ Helper Pseudocode}
中所示算法的直接并行化。
\begin{fcvref}[ln:SMPdesign:SEQ Pseudocode]
\cref{lst:SMPdesign:SEQ Pseudocode}的\Clnref{ifge}必须使用fetch-and-add，
并且局部变量\co{vi}必须在各线程之间共享。
\end{fcvref}
\begin{fcvref}[ln:SMPdesign:SEQ Helper Pseudocode:try]
\cref{lst:SMPdesign:SEQ Helper Pseudocode}的
\Clnref{chk:not:visited,mark:visited}必须
组合成一个CAS循环，CAS失败表示迷宫中存在环。
此清单的\Clnrefrange{recordnext}{next:visited}必须使用
fetch-and-add来仲裁在\co{->visited[]}数组中
记录单元格的并发尝试。
\end{fcvref}

这种方法确实在运行于2.53\,GHz的双CPU Lenovo W500上
提供了显著的加速，如
\cref{fig:SMPdesign:CDF of Solution Times For SEQ and PWQ}所示，
该图展示了两种算法的求解时间的累积分布函数（CDF），
基于500个不同的500乘500随机生成方形迷宫的求解。
CDF在x轴投影上的大量重叠将在
\cref{sec:SMPdesign:Performance Comparison I}中讨论。

有趣的是，顺序求解路径跟踪在并行算法中无需修改即可工作。
然而，这暴露了并行算法的一个显著弱点：
在任何给定时刻，最多只有一个线程可以沿着求解路径取得进展。
这个弱点将在下一节中解决。

\subsection{替代并行迷宫求解器}
\label{sec:SMPdesign:Alternative Parallel Maze Solver}

年轻一代的研究者通常认为从迷宫两头同时开始计算比较好，
这个建议近来在自动求解迷宫问题的文献中
被反复提出~\cite{UMD:CMSC433maze}。
这个方法本质上就是分区，无论在并行操作系统
内核~\cite{Beck85,Inman85}还是在
应用程序~\cite{DavidAPatterson2010TroubleMulticore}层面，
分区一直是一种非常有效的并行化策略。
本节采用这一方法，使用两个子线程从求解路径的
相对两端同时开始，并简要考察其性能和可扩展性后果。

\begin{listing}
\begin{fcvlabel}[ln:SMPdesign:Partitioned Parallel Solver Pseudocode]
\begin{VerbatimL}[commandchars=\\\@\$]
int maze_solve_child(maze *mp, cell *visited, cell sc)	\lnlbl@b$
{
	cell c;
	cell n;
	int vi = 0;

	myvisited = visited; myvi = &vi;		\lnlbl@store:ptr$
	c = visited[vi];				\lnlbl@retrieve$
	do {
		while (!maze_find_any_next_cell(mp, c, &n)) {
			if (visited[++vi].row < 0)
				return 0;
			if (READ_ONCE(mp->done))	\lnlbl@chk:done1$
				return 1;
			c = visited[vi];
		}
		do {
			if (READ_ONCE(mp->done))	\lnlbl@chk:done2$
				return 1;
			c = n;
		} while (maze_find_any_next_cell(mp, c, &n));
		c = visited[vi];
	} while (!READ_ONCE(mp->done));		\lnlbl@chk:done3$
	return 1;
}
\end{VerbatimL}
\end{fcvlabel}
\caption{分区并行求解器伪代码}
\label{lst:SMPdesign:Partitioned Parallel Solver Pseudocode}
\end{listing}

\begin{fcvref}[ln:SMPdesign:Partitioned Parallel Solver Pseudocode]
分区并行算法（PART），如
\cref{lst:SMPdesign:Partitioned Parallel Solver Pseudocode}
（\path{maze_part.c}）所示，
与SEQ类似，但有几个重要区别。
首先，每个子线程有自己的\co{visited}数组，由
父线程传入，如\clnref{b}所示，
该数组必须初始化为全[$-1$, $-1$]。
\Clnref{store:ptr}将指向此数组的指针存储到每线程变量
\co{myvisited}中，以便辅助函数访问，并类似地存储
指向本地访问索引的指针。
其次，父线程代表每个子线程访问第一个单元格，
子线程在\clnref{retrieve}上获取它。
第三，一旦一个子线程定位到另一个子线程已经访问过的单元格，
迷宫就被解决了。
当\co{maze_try_visit_cell()}检测到这一情况时，
它设置迷宫结构中的\co{->done}字段。
第四，每个子线程因此必须定期检查\co{->done}字段，
如\clnref{chk:done1,chk:done2,chk:done3}所示。
\co{READ_ONCE()}原语必须禁止任何可能合并连续加载
或可能重新加载值的编译器优化。
一个C++1x volatile relaxed load就足够了~\cite{RichardSmith2019N4800}。
最后，\co{maze_find_any_next_cell()}函数必须使用
compare-and-swap来标记单元格为已访问，但是
除了线程创建和join提供的排序之外，不需要额外的排序约束。
\end{fcvref}

\begin{listing}
\begin{fcvlabel}[ln:SMPdesign:Partitioned Parallel Helper Pseudocode]
\begin{VerbatimL}[commandchars=\\\@\$]
int maze_try_visit_cell(struct maze *mp, int c, int t,
                        int *n, int d)
{
	cell_t t;
	cell_t *tp;
	int vi;

	if (!maze_cells_connected(mp, c, t))		\lnlbl@chk:conn:b$
		return 0;				\lnlbl@chk:conn:e$
	tp = celladdr(mp, t);
	do {						\lnlbl@loop:b$
		t = READ_ONCE(*tp);
		if (t & VISITED) {			\lnlbl@chk:visited$
			if ((t & TID) != mytid)		\lnlbl@chk:other$
				mp->done = 1;		\lnlbl@located$
			return 0;			\lnlbl@ret:fail$
		}
	} while (!CAS(tp, t, t | VISITED | myid | d));	\lnlbl@loop:e$
	*n = t;						\lnlbl@update:new$
	vi = (*myvi)++;					\lnlbl@update:visited:b$
	myvisited[vi] = t;				\lnlbl@update:visited:e$
	return 1;					\lnlbl@ret:success$
}
\end{VerbatimL}
\end{fcvlabel}
\caption{分区并行辅助函数伪代码}
\label{lst:SMPdesign:Partitioned Parallel Helper Pseudocode}
\end{listing}

\begin{fcvref}[ln:SMPdesign:Partitioned Parallel Helper Pseudocode]
\co{maze_find_any_next_cell()}的伪代码与
\cref{lst:SMPdesign:SEQ Helper Pseudocode}中所示的相同，
但\co{maze_try_visit_cell()}的伪代码不同，
如\cref{lst:SMPdesign:Partitioned Parallel Helper Pseudocode}所示。
\Clnrefrange{chk:conn:b}{chk:conn:e}
检查单元格是否连通，如果不连通则返回失败。
跨越\clnrefrange{loop:b}{loop:e}的循环尝试将
新单元格标记为已访问。
\Clnref{chk:visited}检查它是否已被访问，如果是则
\clnref{ret:fail}返回失败，但只在\clnref{chk:other}
检查我们是否遇到了另一个线程之后，
如果是的话\clnref{located}指示已找到解。
\Clnref{update:new}更新到新单元格，
\clnref{update:visited:b,update:visited:e}更新此线程的已访问
数组，\clnref{ret:success}返回成功。
\end{fcvref}

\begin{figure}
\centering
\resizebox{2.2in}{!}{\includegraphics{SMPdesign/500-ms_seq_fg_part-cdf}}
\caption{SEQ、PWQ和PART求解时间的CDF}
\label{fig:SMPdesign:CDF of Solution Times For SEQ; PWQ; and PART}
\end{figure}

性能测试揭示了一点小小的意外，
如\cref{fig:SMPdesign:CDF of Solution Times For SEQ; PWQ; and PART}所示。
尽管只运行在两个线程上，PART的中位求解时间（17毫秒）
比SEQ（79毫秒）快四倍多。

面对如此戏剧性的性能异常，第一反应是检查
是否有bug，这意味着应该进行严格的验证。
这是下一节的主题。

\subsection{迷宫验证}
\label{sec:SMPdesign:Maze Validation}

大部分验证工作由一致性检查组成，
可以通过在\path{CodeSamples/SMPdesign/maze/*.c}中搜索\co{ABORT()}来定位。
示例检查包括：

\begin{enumerate}
\item	迷宫求解步骤最终落在迷宫之外。
\item	迷宫突然有零行或更少的行或列。
\item	新创建的迷宫包含不可达的单元格。
\item	无解的迷宫。
\item	不连续的迷宫解。
\item	尝试在迷宫外部启动迷宫求解器。
\item	解路径长于迷宫中单元格数量的迷宫。
\item	不同线程的子解互相交叉。
\item	内存分配失败。
\item	系统调用失败。
\end{enumerate}

Paul的妻子进行了额外的手动验证，她非常喜欢解谜题。

然而，如果这个迷宫软件要在生产环境中使用，
无论这意味着什么，明智的做法是构建一个独立的迷宫
\co{fsck}程序。
尽管如此，所有迷宫和解都被证明是完全有效的。
因此下一节更深入地分析了
\cref{sec:SMPdesign:Alternative Parallel Maze Solver}中指出的可扩展性异常。

\subsection{性能比较 I}
\label{sec:SMPdesign:Performance Comparison I}

\begin{figure}
\centering
\resizebox{2.2in}{!}{\includegraphics{SMPdesign/500-ms_seqVfg_part-cdf}}
\caption{SEQ/PWQ和SEQ/PART求解时间比值的CDF}
\label{fig:SMPdesign:CDF of SEQ/PWQ and SEQ/PART Solution-Time Ratios}
\end{figure}

虽然这些算法确实为有效的迷宫找到了有效的解，
但\cref{fig:SMPdesign:CDF of Solution Times For SEQ; PWQ; and PART}
中CDF的图假设数据点是独立的。
事实并非如此：
性能测试随机生成一个迷宫，然后在该迷宫上运行所有求解器。
因此，绘制每个生成迷宫的求解时间比值的CDF是有意义的，
如\cref{fig:SMPdesign:CDF of SEQ/PWQ and SEQ/PART Solution-Time Ratios}所示，
这大大减少了CDF的重叠。
该图揭示了对于某些迷宫，PART
比SEQ快将近\emph{四十}倍。
相比之下，PWQ的速度从不超过SEQ的约两倍。
两个线程就能带来四十倍的性能提升，这当然需要一个解释。
毕竟这不仅仅是易并行的，即可分区性
意味着增加线程不会增加总体计算成本。
而是\emph{\IX{humiliatingly parallel}}（令人震惊的并行）：
增加线程显著降低了总体计算成本，
从而带来了巨大的算法超线性加速。

\begin{figure}
\centering
\resizebox{1.4in}{!}{\includegraphics{SMPdesign/maze_in_way10a}}
\caption{低访问百分比的原因}
\label{fig:SMPdesign:Reason for Small Visit Percentages}
\end{figure}

\begin{figure}
\centering
\resizebox{2.2in}{!}{\includegraphics{SMPdesign/500-pctVms_seq_part-sct}}
\caption{访问百分比与求解时间的相关性}
\label{fig:SMPdesign:Correlation Between Visit Percentage and Solution Time}
\end{figure}

进一步调查显示，
PART有时只访问了不到2\pct\ 的迷宫单元格，
而SEQ和PWQ从未访问少于约9\pct。
关于这种差异的解释见
\cref{fig:SMPdesign:Reason for Small Visit Percentages}。
如果从左上角遍历的线程走到了圆圈处，
另一个线程就根本没有机会访问迷宫的右上部分。
同样，如果另一个线程先走到了方块处，
第一个线程也没有机会访问迷宫的左下部分。
因此PART极有可能只需访问非解路径单元格的一小部分就能找到答案。
简而言之，这种超线性的性能加速是因为一个线程挡住了另一个线程的路。
这简直和数十年来并行编程的经验背道而驰，
我们一直以来都努力保证线程之间\emph{不要}互相妨碍。

\Cref{fig:SMPdesign:Correlation Between Visit Percentage and Solution Time}
证实了在三种方法中，访问的单元格比例与求解时间之间
存在强相关性。
代表PART的散点图斜率小于SEQ，
这说明PART的两个线程访问各自部分的迷宫时
比SEQ的单线程还快很多。
PART的散点图也偏向较小的访问百分比，
说明PART要做的整体工作也较少，这也是
观察到的令人震惊的并行性的原因。
这种令人震惊的并行性还在两个CPU上提供了超过2倍的加速，
如\cpageref{sec:SMPdesign:Problems Maze}中所述。

\begin{figure}
\centering
\resizebox{1.4in}{!}{\includegraphics{SMPdesign/maze_PWQ_vs_PART}}
\caption{PWQ潜在竞争点}
\label{fig:SMPdesign:PWQ Potential Contention Points}
\end{figure}

PWQ访问的单元格比例与SEQ相当。
但是即使访问单元格的比例相同时，
PWQ的求解时间也要比PART长。
这一点的解释见\cref{fig:SMPdesign:PWQ Potential Contention Points}，
图中每个拥有两个以上邻居的单元格上有一个红色圆圈。
每个这样的单元格都可能在PWQ中导致竞争，因为
一个线程进入时可能两个线程正要退出，这会损害性能，
正如本章前面所述。
相比之下，PART只有在找到解时才竞争一次。
当然，SEQ永远不会竞争。

\QuickQuiz{
	既然2D迷宫在两个CPU上实现了4倍加速，
	3D迷宫能否在两个CPU上实现8倍加速？
}\QuickQuizAnswer{
	这是一个很好的问题，留给有足够兴趣
	和勤奋的读者。
}\QuickQuizEnd

\begin{figure}
\centering
\resizebox{2.2in}{!}{\includegraphics{SMPdesign/500-ms_seqVfg_part_seqO3-cdf}}
\caption{编译器优化的效果（-O3）}
\label{fig:SMPdesign:Effect of Compiler Optimization (-O3)}
\end{figure}

虽然PART的加速比令人印象深刻，但我们也不应忽视对顺序实现的优化。
\Cref{fig:SMPdesign:Effect of Compiler Optimization (-O3)}表明，
用-O3编译的SEQ比未优化的PWQ快约两倍，
与未优化的PART性能相当。
如果所有三个算法都用-O3编译，其结果与
\cref{fig:SMPdesign:CDF of SEQ/PWQ and SEQ/PART Solution-Time Ratios}
中所示的类似（尽管都更快），
但是PWQ和SEQ相比并没有快多少，
这与\IXr{Amdahl's Law}~\cite{GeneAmdahl1967AmdahlsLaw}相符。
但是，如果目标是让程序性能比未优化SEQ翻倍，
而不是达到最优性，那么编译器优化是一个值得考虑的选项。

Cache对齐和填充有时通过减少\IX{false sharing}来提高性能。
然而，对于这些迷宫求解算法，在1000x1000迷宫中对迷宫单元格数组进行对齐和填充
\emph{降低}了高达42\pct\ 的性能。
相较于避免false sharing，cache局部性此时更加重要，
对于大型迷宫尤其如此。
对于较小的20乘20或50乘50迷宫，对齐和填充可以为PART
带来高达40\pct\ 的性能提升，
但对于这样小的迷宫，SEQ的性能仍然更好，因为
PART在线程创建和销毁时的开销抵消了它的优势。

简而言之，分区并行迷宫求解器是个有趣的例子，
展示了算法超线性的加速。
如果"算法超线性加速"让你感到理解上的困难，
请继续阅读下一节。

\subsection{替代顺序迷宫求解器}
\label{sec:SMPdesign:Alternative Sequential Maze Solver}

\begin{figure}
\centering
\resizebox{2.2in}{!}{\includegraphics{SMPdesign/500-ms_seqO3V2seqO3_fgO3_partO3-cdf}}
\caption{分区协程}
\label{fig:SMPdesign:Partitioned Coroutines}
\end{figure}

算法超线性加速的存在暗示可以通过协程来模拟并行，
例如，在\cref{lst:SMPdesign:Partitioned Parallel Solver Pseudocode}
的主do-while循环的每次遍历中手动在线程之间切换上下文。
这种上下文切换很直接，因为上下文仅由变量\co{c}和\co{vi}组成：
在实现这种效果的众多方法中，这是上下文切换开销
和访问百分比之间的一个良好权衡。
如\cref{fig:SMPdesign:Partitioned Coroutines}所示，
这个协程算法（COPART）相当有效，
在一个线程上的性能在PART在两个线程上性能的约30\pct\ 以内
（\path{maze_2seq.c}）。

\subsection{性能比较 II}
\label{sec:SMPdesign:Performance Comparison II}

\begin{figure}
\centering
\resizebox{2.2in}{!}{\includegraphics{SMPdesign/500-ms_seqO3VfgO3_partO3-median}}
\caption{不同迷宫大小 vs.\@ SEQ}
\label{fig:SMPdesign:Varying Maze Size vs. SEQ}
\end{figure}

\begin{figure}
\centering
\resizebox{2.2in}{!}{\includegraphics{SMPdesign/500-ms_2seqO3VfgO3_partO3-median}}
\caption{不同迷宫大小 vs.\@ COPART}
\label{fig:SMPdesign:Varying Maze Size vs. COPART}
\end{figure}

\Cref{fig:SMPdesign:Varying Maze Size vs. SEQ,%
fig:SMPdesign:Varying Maze Size vs. COPART}
展示了变化迷宫大小的效果，分别将运行在两个线程上的PWQ和PART
与SEQ或COPART进行比较，带有90\=/percent\-/confidence
误差条。
对于100乘100及更大的迷宫，PART相对于SEQ显示出超线性可扩展性，
相对于COPART显示出适度的可扩展性。
PART在大约200乘200的迷宫大小时超过了相对于COPART的
理论能效平衡点，因为在高频率下功耗大致随频率的平方
增长~\cite{TrevorMudge2000Power}，
因此两个线程上1.4倍的扩展在相同求解速度下消耗的能量
与单线程相同。
相比之下，除非未优化，PWQ相对于SEQ和COPART都显示出
较差的可扩展性：
\Cref{fig:SMPdesign:Varying Maze Size vs. SEQ,%
fig:SMPdesign:Varying Maze Size vs. COPART}
是使用-O3生成的。

\begin{figure}
\centering
\resizebox{2.2in}{!}{\includegraphics{SMPdesign/1000-ms_2seqO3VfgO3_partO3-mean}}
\caption{平均加速比 vs.\@ 线程数，1000x1000迷宫}
\label{fig:SMPdesign:Mean Speedup vs. Number of Threads; 1000x1000 Maze}
\end{figure}

\Cref{fig:SMPdesign:Mean Speedup vs. Number of Threads; 1000x1000 Maze}
展示了PWQ和PART相对于COPART的性能。
对于超过两个线程的PART运行，额外的线程
沿着连接起始和终止单元格的对角线均匀分布启动。
简化的链路状态路由~\cite{BERT-87}被用来
检测超过两个线程的PART运行中的提前终止
（当一个线程同时连接到起点和终点时，标记解已找到）。
PWQ性能相当差，但
PART在两个线程处达到平衡点，在五个线程处再次达到平衡，
在五个线程之后实现了适度的加速。
理论能效平衡点在七个和八个线程的90\=/percent\-/confidence
区间内。
在两个线程处出现峰值的原因是（1）两线程情况下
终止检测的复杂性较低，以及（2）第三个及后续线程
取得有用前进进展的概率较低：
只有前两个线程保证从解路径上开始。
与\cref{fig:SMPdesign:Varying Maze Size vs. COPART}中的结果
相比，这种令人失望的性能是由于运行在2.66\,GHz的
更大更旧的Xeon系统中集成度较低的硬件造成的。

\QuickQuiz{
	为什么把第三个、第四个等线程放在对角线上？
	为什么不把它们均匀分布在整个迷宫中？
}\QuickQuizAnswer{
	确实有许多分配额外线程的方法。
	分配策略的评估留给有足够兴趣
	和勤奋的读者。
}\QuickQuizEnd

\subsection{未来方向和结论}
\label{sec:SMPdesign:Future Directions and Conclusions}

未来还有很多工作要做。
首先，人类在解决迷宫问题时有很多手段，本节只使用了其中一种技术。
其他可能的方法有靠墙行走以排除迷宫的部分区域，
以及根据先前遍历路径的位置选择内部起始点。
其次，对起点和终点的不同选择可能有利于不同的算法。
第三，虽然PART算法前两个线程的出发点很容易选择，
但选择后续线程的出发点有许多种方式。
最优的出发点选择可能与起点和终点的位置有关。
第四，对不可解迷宫和循环迷宫的研究
也可能产生有趣的结果。
第五，轻量级C++11原子操作可能会提升性能。
第六，比较二维迷宫和三维迷宫（或更高维迷宫）算法的加速比
也是一个有趣的问题。
最后，对于迷宫问题来说，令人震惊的并行性指出了
一种使用协程的更高效顺序实现。
那么对任何问题来说，只要出现了令人震惊的并行性，
是否就一定意味着会有更高效的顺序实现？
还是存在本质上令人震惊的并行算法，
其中协程上下文切换开销压过了加速比？

本节展示并分析了迷宫求解算法的并行化。
基于工作队列的常规算法只在禁用编译器优化时性能不错，
这说明之前一些通过高级/高开销语言
获得的结果可能会被优化的进步所推翻。

本节给出了一个清晰的例子，说明从头开始设计并行算法，
而非从已有的顺序实现中衍生并行方案，
反而有助于找出提升顺序实现性能的办法。
在设计阶段就在较高层面考虑并行性，很可能是
一个富有成果的研究领域。
本节将迷宫求解问题从一般性能扩展到令人震惊的并行，
然后又回来了。
希望这一经验能激励人们在设计阶段就开始考虑并行性，
而不是在实现代码后才开始微调程序来实现并行化。

\section{分区、并行和优化}
\label{sec:SMPdesign:Partitioning; Parallelism; and Optimization}
%
\epigraph{知识如果不付诸实践就毫无价值。}
	 {Anton Chekhov}

虽然本章展示了在设计层面就考虑并行性可以得到优异的结果，
但最后这一节想说的是，这还远远不够。
对于搜索类问题（如迷宫求解），
搜索策略比并行设计更重要。
是的，对于特定类型的迷宫，巧妙地应用并行性
发现了一种更高级的搜索策略，但这种幸运
并不能替代对搜索策略本身的明确关注。

正如\cref{sec:intro:Parallel Programming Goals}中所述，
并行化只是众多优化手段中的一种。
一个成功的设计需要关注最重要的优化手段。
虽然我不情愿承认，但那个最重要的优化
可能是也可能不是并行性。

不过，对于很多问题来说并行化确实是正确的选择，
下一章将介绍同步的主力军——锁。
